\section{The flag conformal structure and field-local passage laws}\label{sec:conformal}

The spatial comparison of Section~\ref{sec:cutting} supplies
orientation-preserving homeomorphisms from the discrete flag spheres to the
continuum sphere, with vanishing flag mesh. We must show that the actual
flag uniformizations, expressed in these reference coordinates, converge
to the continuum conformal normalization. A homeomorphism supplied by the
spatial comparison need not itself be conformal.

We first obtain lower bounds for extremal length on the actual flag
surface. These bounds give compactness of its normalized uniformizations
and prevent a continuum from collapsing. They also make each limiting
homeomorphism quasiconformal. The remaining distortion is measured by its
normalized differential tensor. Local conditional laws make this tensor a
field-germ reading; triviality of Gaussian field germs and covariance under
rotations then force it to be the identity. After this conformal
identification, we reconstruct the fields from the area and signed loop
data and remove the auxiliary exploration choices from the profile
kernels. Those field-only kernels are the input to the metric comparison
in Section~\ref{sec:metric}.

\subsection{Extremal length and the local comparison problem}

Use the extremal length supplied by the actual polyhedral flag triangles with their specified
Euclidean geometry. For a family \(\mathcal G\) of curves in a conformal surface, our convention is
\[
 \EL(\mathcal G)=\sup_{\rho\geq0}
 \frac{\bigl(\inf_{\gamma\in\mathcal G}\int_\gamma\rho\,ds\bigr)^2}
 {\int\rho^2\,dA},
\]
where the supremum ranges over measurable densities with positive finite squared integral. Thus
\(\EL\) is extremal \emph{length}, rather than its reciprocal. If a probability on curves has mean
length traffic with density \(\sigma\) relative to \(dA\), then Cauchy--Schwarz gives \(\EL(\mathcal
G)\leq\int\sigma^2\,dA\). Restricting the probability to a subfamily of mass \(p>0\) and
renormalizing multiplies this upper bound by at most \(p^{-2}\). Densities and length traffic of
curve measures can equivalently be expressed in any conformal background metric, using extended
conformal charts at the cone vertices. We use the usual planar extremal-length theorems (Jordan
quadrilateral/cylinder uniformization and the series inequality; bounded upper EL for connections of
two continua of diameters bounded below on the round sphere). In particular opposite
through-direction ELs of a Jordan quadrilateral multiply to one, as do through and essential
closed-curve EL in a Jordan ring, using curves in the open domain with ends allowed on the indicated
sides. In a cylinder, longitudinal paths with uniform transverse parameter give a probability whose
length traffic has squared density norm equal to the through EL.

The reference coordinates used to locate these families are still only
topological. Accordingly, every local test uses the compact constraints and
open enlargements of Section~\ref{sec:kernels}. We compare families by
inclusion or by mandatory subcrossings; we do not assume that topological
approximation alone preserves extremal length.
The conditional spatial kernels retain the embedded input
\(\mathcal X=(\mathcal T,\mathcal E)\): \(\mathcal T\) is the abstract
primary tree, and \(\mathcal E\) is the remaining field and embedding
realization sampled conditionally on it. All spatial product rules below
keep this realization in the conditioning.

\subsection{Local measure classes and aligned shells}
\label{conf:local-classes}

Work for now on ordinary plane patches of the arbitrary-area sphere class. Use the two area-typical
end marks of Section~\ref{sec:continuum}, a central separating choice of full support, and unpinned
gauges (positive densities including for the remaining embedding; end locations can also be
randomized by an independent Möbius change). Require the ends exterior to the bounded buffered patch
under discussion.

On such a patch the real-height law is locally absolutely continuous relative to an ordinary GFF
with free actual-height distribution, and a fixed reference experiment of this class gives
equivalence on strict inner buffers when the ends are outside a disk containing the patch. Indeed in
the cylinder the lateral field is Gaussian and the average process is the BES maximum representation
of Section~\ref{sec:continuum} with randomized translation and constant. On any finite enlarged
interval this latter law is equivalent to Brownian motion with free initial height: the
inside-maximum cases with positive densities and Denisov/Imhof give domination of the whole Brownian
class, and the outside-maximum pieces are absolutely continuous as well.

Local domain comparison in conformal coordinates is just GFF Markov and harmonic cutoff absolute
continuity. \(J\) is the independent ordinary scalar GFF modulo period. Positive changes of global
Palm densities can be undone on inputs. Thus limiting high-probability tests can also use ordinary
local Gaussian field laws with density-change truncation. End/gauge randomization makes the
unrestricted arbitrary-area law class invariant under deterministic similarities and field-constant
changes (with end and field transformation). Restrict back to exterior ends on prototypes.

The next estimate turns high-probability local tests into shells around
every point. It will be used for extremal-length barriers here and for
graph-cost comparisons in Section~\ref{sec:metric}.
Prototype tests take place with margins in a fixed round annular region buffered strictly within an
open annulus \(\Omega=\{p_-<|z|<p_+\}\) about 0. Tests can include local continuum field/metric
regularities, as well as robust profile assertions. Use one arbitrary-area prototype measure class
experiment as above with ends outside a larger fixed disk. Length-perimeter labels use factors
\(2^{-i}\), possibly with a fixed common displacement (\(\tau _i=i\log 2+d\)). Costs use the specified converted units; real fields pulled back
from radius \(r\) about \(z\) to the prototype at this label become
\[
         h(z+r\,\cdot )+2\log r+\tau _i .
\]
Thus scaled continuum perimeter units multiply by \(e^{\tau _i}\), and the discrete perimeter count
unit by \(e^{-\tau _i}\); reference \(D\)-length converts by the metric Weyl rule. One does not
change the microscopic cost on an individual path in such a reinterpretation. Work with joint
limiting extractions and the same kernel versions under all fixed label conversions. Concretely
unconditioned bases at corresponding scales in fixed ratios are reweightings (perimeter mixing
densities), and \eqref{eq:original-15} under converted measurements/calculation histories identifies
the same local readout rules; actual costs may additionally be divided by specified deterministic
units/ratios. Continuum placements can use the scaled Palm class. Thus the probabilities at
different labels need not come from any limiting Weyl- or dilation-invariance of graph observations.

\begin{proposition}[Aligned all-point shells]\label{conf:aligned-shells}

Fix \(\alpha>0\), the buffered annular geometry above, and a
countable common extraction of field and regularized-profile observations
with its fixed label conversions. Let \(E_i\) be a local success test at
label \(i\), with strict band margins, evaluated in the fixed prototype
probability law. For sufficiently small \(q>0\), chosen after these
fixed comparisons, the following implication holds. Suppose that for
each \(N\) in an unbounded set of integers there is
\(G_N\subset[N,2N]\cap\Z\) satisfying
\[
 |G_N|\ge\alpha N,
 \qquad \Prob(E_i)\ge1-q\quad(i\in G_N).
\]
Then, as \(N\to\infty\) through this set, with probability tending to
one every point of any fixed planar compact away from the ends lies,
with central-hole clearance, in at least \(\beta N\) successful shell
holes along its fine-grid ancestry. Here \(\beta>0\) and all the radii
are at most \(e^{-aN}\) for a fixed \(a>0\). For each label, disks of
any fixed large multiple \(P_0\) of these radii can be required to have
bounded overlap.

For each fixed \(N\) only finitely many shells are tested, although the
number of spatial size levels may be exponential in \(N\). The discrete
index is sent to infinity for this finite family before \(N\) is
increased. The same conclusion applies along jointly extracted
subsequences; the required prototype probability is chosen after the
geometry, fraction, density comparisons, and simultaneous-candidate
requirements have been fixed.

\end{proposition}\begin{proof}

\smallskip\noindent\emph{Uniform candidate probabilities and changes of density.}\par

All-order success suffices. Apply the common finite-index and independent-draw conclusions of
Section~\ref{sec:kernels}, in particular Theorem~\ref{kern:finite-vector} and
Corollary~\ref{kern:candidate-draws}. For labels in a common extraction it turns arbitrarily high
prototype probabilities into arbitrarily high local simultaneous-candidate probabilities uniformly
by truncating smallest weights. If considering successive different extractions/probability choices
one can equally include them in a countable common sequence of tests/extractions on the same
prototype inputs: the finite-variant bound is by the common fields/order comparison, using common
chart choices, not by the value of the distance or EL units. For fixed comparable reference types
the functions may be different universal disintegrations; the \emph{alternative} used is indexed by
the same finite choices over the extensions for the whole vector. Joint diagonal countersequence
extraction suffices for uniform probability-loss assertions.

Under Gaussian Dirichlet trials on \(\Omega \) with real-height and \(J\) harmonic corrections
bounded on a slightly smaller region (imaginary constant taken modulo period), this high probability
transfers uniformly with strict buffers: cutoff Cameron–Martin shifts for harmonic data bounded
smoothly there have uniformly bounded finite-moment likelihoods, and fixed restricted Gaussian
domain/background comparisons use absolute continuity as above. Pools remain private.

For any one deterministic ancestry of disjoint trial shells in a limiting experiment, the actual
selected kernels are among the feasible prototype choices at the separately zoomed labels (the input
pushforward is absolutely continuous in that global class, and the same conditional chart rule
applies). Simultaneous-candidate arrays can be sampled as above including the selected observations
by disjoint conditional locality.

To bound a union over such ancestries it suffices to make this comparison for each path separately
conditional on enclosing local fields. Truncate once the field-density change on the enclosing
ordinary patch from a pure log Gaussian experiment; multiply the path probability bounds by that
truncation. Thus there is no density factor per spatial grid point from supplying the sphere fields.
One can for instance restrict to ends outside fixed enclosing neighborhoods, or exhaust that event.
Feasibility and conditional draws in these arguments are only asserted for typical data in the
indicated measure classes.

\proofstage{The grid and the critical change of measure}

Use nested square grids of radii \(r_j=m^{-j}\), \(T=\log m\) with a large fixed integer \(m\), cell
side \(c_*' r_j\) very fine relative to the hole. Around the cell center \(z_j\) let \(H_j\) be the
real average on radius \(R r_j\), with \(R\gg P_0+p_+\) fixed. Ancestor shells are disjoint and
these circles are off the shells on that ancestry, increasing \(m\) as needed. Use
\[
            X_j=2Tj-H_j .
\]
Take \(Tj_0\sim aN\) for very small fixed \(a>0\), and \(J_N=\lceil\exp (A N)\rceil \) as last
index, for a sufficiently large constant \(A\). In the enclosing pure Gaussian calculation use
covariance \(-\log |u-v|+\mathrm{const}\) on a containing disk for the real field, and the same
modulo constants plus a phase for \(J\). With probability tending to one all starting \(X_{j_0}\)'s
are below the target range with any desired fixed clearance (Gaussian grid union bound, \(4a<\log
2\)). Along each path stop at the first \(K\) with \(X_K\ge 3N\log 2\).

For a fixed grid path to \(k\), tilt by \(\exp (2H_k-2\operatorname{Var}H_k)\). By nested-circle
harmonic averaging the increments of \(X_j\) to \(k\) now have independent centered Gaussian laws of
variance \(T\); \(X_{j_0}\) has bounded mean and variance \(O(N)\). The number of grid cells at level \(k\) is \(O(e^{2Tk})\).
Since \(\Var(H_k)=Tk+O(1)\), multiplication by the inverse tilt gives
\[
 O(e^{2Tk})\exp\{-2H_k+2\Var(H_k)\}
       =O(e^{2X_k}).
\]
Thus the grid cardinality is absorbed into the terminal value of the
centered tilted walk.

Paths not stopped by \(k=J_N\) are negligible after this union: under tilt the probability of
staying below the barrier is at most \(C(N+\sqrt {\log J_N})/\sqrt {J_N}\), by reflection for the
Brownian interpolation allowing the maximum intersample oscillation (add a negligible tail), and
averaging the start. Choose \(A\) accordingly.

For first-stop paths to \(k\le J_N\) with start good, \(e^{2X_k}\le \exp (6N\log 2+2|\Delta X_k|)\)
with \(\Delta X_j=X_j-X_{j-1}\). Thus tilted bad-on-stopping probabilities \(\exp (-B N)\) at
arbitrary large required \(B\) suffice by Cauchy–Schwarz, including the sum over \(k\).

\smallskip\noindent\emph{Bounded overshoots and clumping.}\par

Assign tested labels to first sampled hits \(j\) of \(X_j\ge \tau _i\). Clump those at the same
\(j\). Except at the above arbitrarily small exponential probability, a positive fraction of labels
occur with last increment \(|\Delta X_j|\le U\) for fixed large \(U\); then each such clump has
bounded size. Indeed one can continue the tilted Brownian walk beyond \(k\) in estimates. At the
first passage of a next level from below, the last increment has a uniform sub-Gaussian upper tail.
To see this at cutoff \(u\) first wait if necessary for the continuous interpolation by Brownian
motion to hit the level minus \(u/3\). Immediate crossing by the next sample then costs \(\exp (-c
u^2)\). Otherwise wait up to \(\exp (c_0u^2)\) samples, bounding large increments by a union with
small \(c_0>0\), and longer delay by reflection and intersample oscillations as before. This works
conditionally when proceeding from clump to clump.

Labels lost at a large-increment clump number at most \(1+|\Delta X_j|/\log 2\); conditional
exponential moments of this loss restricted to \(|\Delta X_j|>U\) tend to one at each fixed
exponent. Iteration proves the estimate. Keep a representative per remaining clump.

\smallskip\noindent\emph{Harmonic data at the selected levels.}\par

The harmonic data off the disjoint shells on this path have bounded needed norms for most
representatives, and differences of \(H_j\) with the corresponding same-radius averages centered at
all grid neighbors within \((3P_0+1)r_j\) are bounded there as well, with arbitrarily small
exponential failure.

Here are details to account for random hit times. Measure harmonic profiles minus own outer circle
values in \(L^2\) on a slightly inset annulus enough for all cutoff comparisons, and include the
neighbor differences. These Gaussians are zero-mass boundary charges with uniformly bounded kernel
variances at each scale (annular Poisson densities and uniform circle charges). Across deterministic
levels their cross covariances decay geometrically by fine-charge cancellation and separation. Thus
covariance on any level sublist has bounded operator norm and trace \(O(\#\mathrm{levels})\), also
after conditioning on the \(X\) path. Imaginary profiles use an arbitrary lift with constant removed
here.

Conditional real means at \(j>j_0\), including under the terminal tilt, are bounded in profile norm by
\[
 C\Big (1+|\Delta X_j|+\sum _{l=j+1}^k m^{-(l-j-1)}|\Delta X_l|\Big ).
\]
Indeed covariance with coarser path-circle values vanishes exactly; with finer ones successive
differences decay by smoothness away from the charges. Use independent circle increments/regression.
At the selected clumps with previous increment bounded these squared-mean bounds sum to \(O_B(N)\)
off an \(\exp (-B N)\) failure. For each offset \(s\ge 0\), increments one step after the distinct
hit indices plus \(s\) satisfy a sum-of-squares bound \(C_B(s+1)^2N\) simultaneously with this
error, by conditional Gaussian exponential moments at those sequential stopping times (use the
extended walk and all clump leaders).

Now condition on the path; the covariance bound and Gaussian squared-norm exponential estimate
finish the profile assertion, by taking norm cutoffs large and discarding just a small fraction.
Harmonic interiors then have the requisite smooth bounds, with real constant aligned up to
\(U+O(1)\).

\smallskip\noindent\emph{Delayed hits, independent trials, and overlap.}\par

Let \(W\) exceed the retained neighbor-error bound with slack. We can further keep hits \(j\) for
which all \(X_{j'}<\tau _i-W\) for \(j_0\le j'<j-L\), for large fixed \(L\), still losing only a
small fraction. The delay from first lower-barrier hit to target hit exceeds \(L\) with conditional
probability \(O((W+\sqrt {\log L})/\sqrt L)\). Split labels into finitely many spacing classes
greater than \(W\), giving sequentially bounded delay failures and hence again arbitrary exponential
estimates.

Condition for actual trials off the shells (including the lifted exterior data); the stopping,
harmonic and label-hit data needed to run trials are exterior, and the tilt changes no Dirichlet
interiors. Neighbor errors if needed can just be discarded by the preceding bounds separately
without conditioning interior trials on them. Conditional independence and the uniformly high
all-candidate trial probabilities now lose only a small fraction of the many representative hits,
again except for arbitrary exponential failure. This proves the union estimate and all-point
assertion.

If enlarged successful disks for one label intersect at indices \(j<l\) with \(l-j>L\), the ancestor
center of the fine one at \(j\) is a tested neighbor of the coarse one, contradicting the
lower-delay requirement. This gives bounded overlap per label. In later single-shell Weyl changes
supported within \(P_0 r_j\), the own and these neighbor outer circle values and all coarser delay
checks on that ancestry are unaffected; the overlap criterion for eligible outputs therefore still
holds.

\end{proof}

\subsection{Prototype annuli and extremal-length quantiles}
\label{conf:el-setup}

Here are details on fixing a reference experiment. Use \(P_j=2^j\) as perimeter units in
unconditioned pair comparisons, with the corresponding path-time and volume units. Mix perimeters by
smooth positive densities. Continuum inputs may be weighted to any equivalent finite probability,
using the arbitrary-area class, full-support gauges (with two area-typical ends, also allowing
independent Möbius variation), and when using ends a central loop separating them with a positive
counting density. This has the pair class by the Palm rule (end choices can be added on opposite
sides).

Keep \(O\) with the convention that a blue-side field mean is 0 modulo 4 in height-gap units (here
and below meaning units of \(\lambda\)); the fair parity and signed increments above give this
coupling. A loop sign can be specified by its actual orientation, fixing which direction gives an
increasing height jump toward which side in the level coupling. Reversing which side is the caller
then reverses both the relative clockwise convention in the disk and the signed increment toward the
positive offspring. Thus these are the preassigned fair ring orientations (central orientation fair
too), not a new coloring sample at each call. A global reversal of the height convention makes no
difference.

For the next quantiles take a round annular band compactly in the planar part and keep margins on
either side; condition the prototype ends to be outside a larger fixed disk as in
Proposition~\ref{conf:aligned-shells}. One precise finite comparison convention couples the
countable data to the true limiting \({\cal T}\) with vanishing errors in the compact tests above,
then samples the placements/fields (including end/gauge variables) given the limiting input
independently of extra readouts. Continuum input weights or restrictions can be imposed on this
experiment afterwards, with ordinary truncation if needed. On any limiting array extraction the
resulting coupling is the same primary disintegration described above: the matched reference input
has a fixed law, so conditional placement expectations evaluated on it can use Lusin approximation
under that law. Along further good subsequences the small-error comparison homeomorphisms thus
exist.

Choose a discrete Jordan annulus following the band boundaries with vanishing error in probability
there (for example by pulling back using the approximate homeomorphisms and taking countable smooth
approximations with the same comparison margins). It separates inner from outer as indicated; a
temporary comparison failure can be assigned value one instead. Choices need not be sharp-boundary
local. The following two transfers use only strict buffered readings:

\begin{itemize}

\item a large through EL gives a barrier for true hole-to-exterior traversal of a slightly larger
band, eventually testable with flag covers;

\item a small through EL gives unit flows confined to an ordinary compact in the band enlargement
with fixed positive advance and squared length-density norm at most that small value.

\end{itemize}

Indeed one can put intermediate families between wide-band mandatory subtraversals and the measured
band (so all intermediate tested paths contain measured traversals); these use opposite ports,
confinement and stars in the annular buffer by the all-lift comparison. An EL bound from below gives
charging densities there and they need only be used within that buffer. The second assertion uses
cylinder traffic in the measured annulus itself. Its resulting strict interior port tests also
transfer by intermediate covers. All these statements use family containment or mandatory subcurves,
not stability of extremal length under topological approximation alone.

Fix small \(\epsilon>0\) and let \(b_j=b_j(\epsilon)\) be a positive lower \(\epsilon\)-quantile of
the measured EL (chosen monotonically in \(\epsilon\)) with
\[
 \mathbb P(\mathrm{EL}\ge b_j/2)\ \ge 1-2\epsilon,\qquad
 \mathbb P(\mathrm{EL}\le 2 b_j)\ \ge\epsilon/2.
\]
The measurements are finite strictly positive. High probability in the first inequality gives, in
limiting extractions, high probability of the buffered eventual lower barrier with a weakened
factor, in units of \(b_j\) or any smaller unit. For instance take a coupling with convergent
triangle-cover variational costs in the smaller units; if large measurements occur on a subsequence,
between the two band margins finite upper-image covers give intermediate tests with large cost on
that subsequence and hence eventually a slightly diminished bound. Comparison failures of
nonvanishing spatial error can be avoided by further subsequencing.

This works also at fixed backward offsets \(i\), using experiments at \(P_{j-i}\) and localized
profile laws of the common extraction with the unit conversion of \eqref{eq:original-15}. Concretely
the unconditioned discrete bases are smooth lattice reweightings in the pair class, and in the
fields, adding \(i\log 2\) scales perimeter by \(2^i\) and area by \(2^{2i}\). Similarity changes
with the charge term only change the embedding. The embedded instructions, elapsed time/jump
labeling in converted units, and the conditional chart formulas consequently agree for the
correspondingly transformed inputs (with joint subextraction of reference experiments);
placement-density changes if needed use their supported Palm laws before reading probabilities. The
same microscopic EL has no multiplier. Thus one uses the high probabilities of the fixed prototype
input distribution at each converted label, with actual selected kernels among its alternatives when
separately zoomed there, as in Proposition~\ref{conf:aligned-shells}.

\subsection{Germ determination}
\label{conf:germs}

\begin{lemma}[Deterministic germ readings]\label{conf:germ-determinism}
Consider a jointly extracted array of local robust profile readings, in any of the full-gauge
arbitrary-area comparison classes, or in the area-one class with an unpinned embedding. Let its germ
variables be bounded and spatially measurable, and defined by countably many robust family or array
tests in arbitrarily small neighborhoods. Tests with variable centers are read through intermediate
countable margins. At Lebesgue-almost every deterministic point their conditional laws given \(h,J\)
depend only on the germs of these fields. For Lebesgue-almost every pair of distinct deterministic
points these conditional laws multiply. Each such variable consequently has a deterministic value
almost everywhere in space. Probability-zero and probability-one germ assertions are transported by
rotations about the center.
\end{lemma}
\begin{proof}

Almost surely \(x\) avoids ends and the chosen central seam and has arbitrarily small whole \(O\)
loops about it with Jordan interiors \(D\) containing it. In any larger field buffer choose such a
loop entirely inside with room. It is shared, together with both-side labels, in two conditional
extensions agreeing on that field buffer by wall locality; given two extensions one may choose it
small enough to avoid both central choices (and have its interior toward \(x\) in both rooted
hierarchies). The signed/colored system uses \(J\) as above. This disk is a complete positive birth,
with a new canonical source and an actual fresh subtree.

Strictly interior readings there have the disk kernel, given the disk's path/volume input, its
perimeter, and its protected placements, irrespective of caller-phase, ancestors and the boundary
point where the loop was found. To see this in \eqref{eq:original-15}, chart through that designated
birth, retain the initial unexamined internal filling, and use only tests protected in its interior;
no exterior slot tests are needed there. Its size \(l\) is a full free lattice coordinate (outer
total on this designated face is free, \(l\) equals that total plus the discrepancy); there is no
projected lattice restriction or residue left. Thus the reference rule here is just smooth size
mixing followed by the fresh colored \(T_l\) and its compatible source/tree rule, and disintegration
on the interior subtree data and size.

Reroot the slots at its first fresh uniform source before calculating that rule. Cyclic rerooting
preserves the entire conditional map law at exact size (boundary edge \emph{occurrences} with fixed
boundary color), and the remaining fresh cyclic offset conveys no shape information. In particular
where the parent matched a circumference point is merely a label rotation when expressing the
all-target and embedded internal data from the source. This argument uses no shape-independence from
the field after disintegration.

At fixed volume reserve the nuisance disk away from the protected birth, as in the finite-cylinder
proof of \eqref{eq:original-15} in Section~\ref{sec:kernels}. Choices of birth charts can use
successes on open intervals and positive-support cuts; the versions hold for all actual interior
disks needed here.

Placements and length parametrization, full canonical hierarchy and level traversals within the disk
are local to the fields containing it, using its own source family. Integrate the sources there by
their actual fresh independent conditional law from the continuum coupling. Thus two conditional
extensions can share these data in sufficiently small common \(D\)'s and have identical integrated
kernels for germ readings; end choices and exterior sources do not change the result.

Two distinct test points use disjoint such disks and the product rule, including independence of
their internal source choices. This proves the claims (one may compare conditional iid full-input
copies for each fixed larger buffer and then shrink the buffer; simultaneous product uses full
inputs first).

By spatial measurability, integrating bounded functions of the variables on rational balls now has
zero variance conditional on the fields. Differentiation shows the variables themselves agree almost
surely at area-almost every point with their field-conditional versions (area here is ordinary
planar area).

Ordinary GFF germs are trivial at such deterministic points, also for the independent pair of fields
here. For a Dirichlet field the decreasing Gaussian spaces shrink trivially by density of Dirichlet
Cameron–Martin functions vanishing near the point (zero point capacity), hence the same for
generated sigma fields by Gaussian chaos. For the density statement, use logarithmic cutoffs
which vanish in a disk of radius \(\varepsilon^2\), equal one outside the disk of radius
\(\varepsilon\), and have Dirichlet energy \(O(1/\log(1/\varepsilon))\). Multiplying smooth
Cameron--Martin functions by these cutoffs and then smoothing proves the required approximation.

Modulo-period field restrictions locally are covered by ordinary actual lifts. Actual real-height
restrictions off ends use Gaussian absolute continuity as above, also at area one by
Lemma~\ref{kern:area-one-transfer}. If needed restrict to events of ends outside deterministic
neighborhoods with probabilities tending to one: under domination on one such neighborhood
triviality of the inner completed tail follows by conditional-density martingales, and then by
exhaustion.

Hence at almost every fixed \(x\) such a variable is deterministic. Probability-one or zero statuses
transport under rotations about \(x\). Indeed the full input classes with free embedding have
equivalence under those similarities, and all readings are intrinsic from the primary cover arrays,
with conditional placements from the continuum law; under embedding rotation the arrays and path
input are unchanged and the spatial tests rotate exactly (or via the same strict margins). This
applies before any claimed invariance of passage profiles, and one can unrestrict any fixed
prototype cutoff first. Countably many rotations suffice.

\end{proof}

\subsection{Nondegeneracy of the extremal-length scale}

We now rule out a vanishing prototype EL unit. Such a unit would give
flows of very small energy in the actual flag metric. The shell barriers
convert their traffic into a nonzero, absolutely continuous Euclidean
arclength measure. Rectifiable paths then give nearly straight local
passages on a set of positive planar area. Germ determination transports
these passages to a transverse direction, forcing both crossing
extremal lengths of one quadrilateral to vanish, contrary to reciprocity.
The proof keeps the spatial buffers throughout this argument.


\begin{proposition}\label{conf:el-nondegeneracy}
For all sufficiently small fixed \(\epsilon>0\), the quantiles defined in Section~\ref{conf:el-setup} satisfy

\begin{equation}\label{eq:original-16}
\liminf_{j\to\infty} b_j(\epsilon)>0.
\end{equation}
\end{proposition}
\begin{proof}
\smallskip\noindent\emph{A countersequence and one common finite index.}\par

If false for a positive quantile, there are failures down a sequence of quantile probabilities
tending to zero (one can take monotone lower quantiles). For each extract at new lows \(j\to\infty,\
b_j\to0\). Then \(b_{j-i}\ge b_j\) for fixed positive offsets eventually. Work in EL units
\(b=b_j\). Choose one of these probability levels sufficiently small for the high-probability shell
test with a fixed band, slack and alignment constants, taking joint diagonal subextractions.

Here the all-candidate probability loss can indeed be taken uniformly: put the countably many
vectors of converted prototype conditional kernels on the same full prototype fields and source
choices. Finiteness over conditional extensions with pools in
\eqref{eq:original-14}--\eqref{eq:original-15} was by common local chart/order comparisons for the
whole countable list, not with a variation bound depending on EL units or on the individual
subsequence likelihood of a budget value. Truncate smallest positive weights of alternatives once
for this list. Field density and harmonic cutoff changes in Proposition~\ref{conf:aligned-shells}
are fixed as well. The remaining upper (small-EL) witness probability at the chosen quantile is
strictly positive, not needed uniformly.

\smallskip\noindent\emph{Charging densities from the aligned shells.}\par

Work on ordinary planar compacts in this last extraction. With probability tending to one as
\(N\to\infty\), good holes at aligned labels in \([N,2N]\) cover each point with multiplicity at
least \(\beta N\), their radii \(r\le e^{-aN}\), with margins and bounded overlap of enlarged disks
per label. This is Proposition~\ref{conf:aligned-shells} with the eventual lower barrier and all
offsets available. At each such annulus, for the discrete index sufficiently large, there is a
length density supported inside its buffered disk charging hole-to-exterior traversal by \(r\), with
squared density norm at most \(C r^2/b\).

For a finite open union \(U\) of ordinary rational boxes in the compact region of interest, retain
the holes needed over \(U\), and sum these densities divided by \(N\) to get \(g\). For each fixed
\(N\) passing, eventually
\[
 b\int g^2\ \le\ C\,|U^{+C e^{-aN}}|,
\]
where the right uses Euclidean area and enlargement. Indeed per label the sum of squared radii is
bounded by packing there and density supports have bounded overlap (one can also use the uniform
small-error homeomorphisms and an enlarged buffer factor). Use Cauchy–Schwarz over labels, not over
the potentially much larger number of Euclidean size levels.

Every path piece in projected \(U\) of diameter \(l\) has \(g\)-length at least \(c l-O(e^{-aN})\).
For \(l\) much larger than the maximal radius, each used hole meeting the path with central
clearance thus forces a traversal. On a coordinate projection interval covered by the path integrate
the central-hole multiplicities at path points; each such interval contribution from one hole is at
most a constant times its radius. All comparisons persist for all lifts eventually at this fixed
\(N\) by strict slack. Charges add on disjoint path subintervals. Use countably many rational tests
and an almost surely passing sequence of \(N\)'s by diagonalization.

\smallskip\noindent\emph{Compactness and arclength of small-energy flows.}\par

Keep the event on which small measurements occur along a further subsequence with the successful
compact comparisons; it has positive probability by the quantile inequality in the prototype
coupling and the almost sure eventual budget statements. All conclusions about limiting local data
can now be proved pathwise there. Use the above unit flows, with mean length traffic of squared norm
at most \(2b\), confined to a fixed ordinary compact and making fixed positive advance in
projection.

By Cauchy–Schwarz with the \(g\)-bound, for each fixed \(s>0\) the expected number of disjoint
advances by at least \(s\) is bounded uniformly eventually (take \(N\) sufficiently large depending
on \(s\)). Thus projected path probabilities can be extracted, up to increasing time changes, on
continuous paths.

For detail, mark each first move of size \(2^{-k}\) from the preceding mark at that resolution,
beginning from the start. Eventually the counts are tight at any finite collection of resolutions.
At resolution \(k\), distribute clock mass \(c\,2^{-k}\) equally among the
successive first-exit intervals, continuously over each interval, and add a
strictly increasing continuous base clock of unit mass. Normalize the total
clock. An advance larger than a fixed multiple of \(2^{-k}\) contains a
complete first-exit interval at the next finer resolution. On an event with
bounded exit count at that resolution, it therefore consumes a uniformly
positive amount of clock time. Tightness of these counts gives the required
fine-scale modulus in probability and continuous parametrizations.

Keep a limiting probability with positive advance.

For each finite union \(U\) as above its expected arclength in \(U\) is at most \(C\sqrt{|U|}\).

Indeed take sums of projected advances on finitely many disjoint intervals confined compactly inside
\(U\); with their number bounded, the same bound with enlargement and additive errors \(O(e^{-aN})\)
times the number follows for the expected supremum from \(g\) and the traffic norm. Strict lower
tests on those sums pass under uniform convergence; then increase \(N\), and exhaust interval
numbers and partitions. In particular paths are rectifiable almost surely, and the integrated
arclength measure on this band enlargement is absolutely continuous and nonzero.

\smallskip\noindent\emph{Transverse tangent families.}\par

Ask for path families at every sufficiently small dyadic radius \(s\) about \(x\), crossing from
projection less than \(-s\) to more than \(s\) in some fixed unit axis about the center, with
transverse deviation \(<s/4\), longitudinal extent within \((-3s,3s)\). Allow reversal and require
strict clearance, witnessed with intermediate covers forcing such passages whose limiting EL upper
bound is finite in units \(b\) (possibly radius-dependent). This is a germ reading. Use a countable
fine net of axes. For arclength-almost every interior point of almost every path from the limiting
flow there is a straight nonzero arclength tangent, hence for a sufficiently good axis all these
requirements on the path hold with extra slack.

Paths in the support realizing strict passage give positive mass of paths making the required
subtransit at each such fixed radius. This transfers to the finite flows, hence finite normalized EL
by restriction to that mass and the traffic bound (any charging density applied to those subtransits
charges at least as much along the full paths). Choose covers intermediate within the slack so these
bounds register along the extraction.

The integrated arclength is absolutely continuous and nonzero. Hence,
on the positive-probability flow event, some axis in the countable
net has its germ assertion on a set of positive planar area. Countability
therefore supplies one fixed axis for which this happens with positive
probability. Fubini gives
a positive-measure set of deterministic points where that assertion has
positive probability. Lemma~\ref{conf:germ-determinism} makes the assertion
probability one at almost every such point; its rotation conclusion does
the same for the perpendicular axis. At one common sufficiently small fixed
radius the two tested families both force crossings of a common topological square in transverse
directions with clearance, by the homeomorphism comparison. The two quadrilateral ELs are each
bounded above by \(O(b)\) eventually by mandatory subcrossings, contradicting product one. This
proves \eqref{eq:original-16}.

\end{proof}

\begin{corollary}[Raw extremal-length budgets]\label{conf:raw-el-budgets}
In every diverging perimeter extraction, the following conclusions hold along the good
spatial-comparison subsequences.
\begin{enumerate}
\item The family of all curves making a fixed positive spherical advance has extremal length
bounded below by a strictly positive constant, eventually along the extraction.
\item For circular rings in an ordinary planar chart, with a fixed
comparable-radius ratio and fixed comparable buffer margins, there is a
strictly positive lower bound on through extremal length that is uniform
in radius and location. For each fixed ring test the bound holds
eventually along the extraction; the eventual discrete index may depend
on the ring.
\end{enumerate}
\end{corollary}
\begin{proof}
In raw EL units ordinary planar compact tests in arbitrary diverging perimeter extractions enjoy the
same eventual density budgets with \(b\) replaced by 1 and fixed positive constants. For powers of
two use \eqref{eq:original-16}, taking the quantile probability sufficiently small for the common
candidate weights and shell tests. For general units compare simultaneously to nearest powers of two
and their fixed offsets by \eqref{eq:original-15}, after extracting the bounded ratio (convergent
unit conversions, using the limiting shift).

All constants in the resulting advance and norm inequalities can be fixed across ordinary compact
locations and countable circular-ring tests in that planar chart; local density truncations there
only affect the success probabilities tending to one. These are localized almost sure limiting
profile assertions, using intermediate compact/open margins where required. Thus they hold also at
exact unit sphere volume via \eqref{eq:original-15} and the area-one field matching, on compacta off
marked ends before three-point normalization. The point at infinity can equivalently use the
inverted chart.

For the first assertion, every curve making the prescribed advance makes a fixed compactly
confined advance in one of finitely many ordinary tests avoiding the isolated exceptional
points. Sum the corresponding densities. Countable compact exhaustions also cover random end
locations.

For the second assertion, use traversals in a domain of area \(O(r^2)\) at radius \(r\), and take
\(N\) sufficiently large for each fixed test. A radius ratio of two suffices, and comparable
buffered versions follow from the same estimate. It is enough that the traversed subbands avoid
the ends, by exhausting them compactly within the punctured chart. Thus the estimate also covers
shrinking rings centered at a missing point. Countably many strict tests give the conclusion
for every center and radius with comparable margins. For any fixed finite
collection of rings, take the largest of their eventual indices. This is
the order used below when stacking finitely many rings and then increasing
the number of rings.
\end{proof}

\subsection{Identification of the uniformizations}

\begin{proposition}\label{conf:flag-uniformization}
Include three fresh independent samples from normalized area in the spatial comparison. The actual
orientation-preserving uniformizations of the flag surfaces, normalized at these samples, converge
along every good extracted subsequence to the prescribed conformal normalization of the continuum
sphere. The area, full nested-loop, and maximal flag-mesh comparisons therefore hold in these actual
uniformization coordinates, in their common joint law.
\end{proposition}
\begin{proof}
\smallskip\noindent\emph{Equicontinuity and exclusion of collapsed fibers.}\par

Corollary~\ref{conf:raw-el-budgets} concerns actual discrete flag-surface extremal length. Include
three fresh independent area marks in the comparison (use the normalized measures), which tend to
three distinct continuum area points by the area coupling. Write \(f_n\) for their normalized flag
uniformization expressed by the projection homeomorphism in reference spherical coordinates. These
homeomorphisms are equicontinuous along the good subsequences in the limit: at any hypothetical
accumulation location of large oscillations on shrinking scales stack disjoint comparable-radius
circular rings toward it using Corollary~\ref{conf:raw-el-budgets}(2), making the through EL
arbitrarily large by the series inequality. Yet the inner and outer complementary closed continua in
image then have diameters bounded below, the latter containing at least two separated normalized
marks for small outer radius, contradicting the planar connection bound. Thus subextract to
continuous surjective \(f\). Each fiber is connected by uniform convergence of sphere homeomorphisms
(use inverse images of small closed balls).

No nontrivial such continuum can collapse: take shrinking neighborhoods of it and inside them simple
arcs of uniformly positive advance with Jordan disk thickenings, possible by connectedness. On the
other side fix a nontrivial closed disk with image separated. The connection EL between the two
disks in the discrete image is bounded above uniformly by annular duality, since all essential ring
curves in projection make a fixed positive advance (otherwise both disks reach outside one ball
containing the curve, excluding essential winding), and Corollary~\ref{conf:raw-el-budgets}(1) is
uniform for them independently of the thinnings.

Collapse instead forces this connection EL to infinity by uniform convergence, taking the limits
successively. Hence \(f\) is a homeomorphism.

\smallskip\noindent\emph{Quasiconformality.}\par

On ordinary compacts \(f\) is quasiconformal. Concentric rings in its image (meaning images by \(f\)
of concentric test rings) with radius ratio, say, 4 in the source have the uniform positive through
EL lower bound by using nested subbands for \(f_n\) and uniform convergence. Therefore locally in
planar charts
\[
 \max_{|y-x|\le r}|f(y)-f(x)|\ \le K\min_{|y-x|=4r}|f(y)-f(x)|.
\]
Indeed a diverging ratio gives arbitrarily many radial log-scales around the image center with both
complementary continua hitting each circle, bounding the connection EL to zero by integrating along
those circles. Put
\[
 H_{1/4}(x,f)=\limsup_{r\downarrow0}
 \frac{\max_{|y-x|=r/4}|f(y)-f(x)|}
      {\min_{|y-x|=r}|f(y)-f(x)|}.
\]
The preceding estimate gives \(H_{1/4}(x,f)\le K\) at every
point of an ordinary compact patch. Applying
\cite[Theorem~1, with $\alpha=1/4$ and any $H>K$]{CristeaQC}
on an open disk compactly contained in that patch proves local
quasiconformality. In particular \(f\in W^{1,2}_{\rm loc}\), and its
Jacobian is positive almost everywhere
\cite[Theorem~3.1.2]{AIMQuasiconformal}. The normalized differential
tensor is therefore defined almost everywhere off the marks.

\smallskip\noindent\emph{The normalized differential tensor.}\par

Finally \( (Df)^t Df/|\det Df|\) in local conformal coordinates at its almost everywhere
nondegenerate differential is itself a germ reading (take bounded injective coordinates for it).
Indeed through moduli of every fixed small rectangle in its image can be read by varying strict
endpoint ports and transverse confinement. Intermediate families encompassing through paths of the
rectangle and forcing subtransits of a slightly wider-transverse shorter-longitudinal rectangle
squeeze costs from below as the margins shrink. Explicitly, allow confinement up to a small outward
error, with the two ports neighborhoods of the complete two end sides. Choose cover tests enveloping
the exact quadrilateral passages yet still inside these port/confinement allowances by the all-lift
comparison. Their EL is bounded above by the through EL for the rectangle, and below by the through
EL in the wider-shorter one lying with transverse/end clearance (every allowed path pays for a
subcrossing). In these inequalities the exact rectangles at finite stages use the spatial
homeomorphisms; their conformal images use the actual flag uniformization. Uniform convergence of
the homeomorphisms and inverses gives the Jordan quadrilateral modulus limits used here (classical
continuity under uniform convergence of parametrized Jordan boundaries and sides, or use further
strictly nested crossing envelopes). Thus the values can equally be taken from the local variational
arrays, regardless of other information used to extract \(f\).

Shrinking the rectangles reads the corresponding moduli under the differential by the same
continuity. For example for side vectors \(Lv,w\) along orthogonal unit vectors, the through EL in
the long direction after an invertible linear map \(A\), divided by \(L\), tends to \(|Av|^2/|\det
A|\): compare uniform parallel lines and the constant area density. Countable axes and aspect ratios
recover the tensor.

Apply Lemma~\ref{conf:germ-determinism} to bounded injective coordinates
of this tensor. At almost every deterministic point \(x\), its value is
a deterministic positive symmetric matrix \(T_x\). Rotating the source
coordinate about \(x\) by a quarter turn \(R\) rotates every buffered
rectangle test and changes its tensor to \(R^tT_xR\). The input law after
this rotation is equivalent to the original free-embedding law, and the
readings transform by exactly these rotated tests. The probability-one
assertion that the tensor equals \(T_x\) must therefore also hold for
\(R^tT_xR\). Consequently
\[
                         R^tT_xR=T_x.
\]
A symmetric matrix commuting with a quarter turn is scalar. Since
\(\det T_x=1\), it is the identity.  Thus \(f\) is conformal away from the
finitely many marks by the distortion-one case of the analytic criterion
\cite[Theorem~2.5.4 and the following paragraph]{AIMQuasiconformal}.
Only now do we remove those points: continuity of the sphere homeomorphism
makes it bounded near each omitted point in a target coordinate chart, so
the holomorphic removable-singularity theorem extends it conformally there.
It is therefore the prescribed orientation-preserving M\"obius normalization.

The argument applies also in the unrestricted area classes (normalizing by area probability samples).

\smallskip\noindent\emph{Transfer of the joint observations.}\par

In particular the topological comparisons for area, the full nested loops (also with their auxiliary
signs and colors), and maximal flag mesh now hold under the \emph{actual} flag uniformizations with
common joint laws as identified before, since every such uniformly convergent subsequence has this
normalization. Removing central bias on the fixed finite sphere by \eqref{eq:original-3} and the
count exhaustion gives the demanded area and CLE law with independence in the surface sense and
vanishing mesh there. No metric convergence is used for this step.

\end{proof}

\subsection{Reconstruction of the fields and removal of source choices}
\label{conf:source-setup}

We now use Proposition~\ref{conf:flag-uniformization} to obtain field-local laws for the regularized
graph profiles. Work with limiting \emph{regularized} graph profiles in true marked-uniformization
coordinates followed by an independent full-support Möbius gauge, in the arbitrary-area pair/sphere
experiments; keep unnormalized area in the comparison. The gauge convention is equivalent in measure
class to the previous free-embedding one (conditional placement-density changes are allowed). Actual
profile tests with spatial margins can now be kept directly in those coordinates at finite stages
and have the same strict/regularized readings as the previous \({\cal T}\)-based ones by uniform
conformal comparison.

The three normalization samples for this purpose can be added last in the \({\cal T}\) coupling, by
the conditional area convergence, and their matched Möbius coordinate changes applied to the earlier
arrays. Given full inputs including those marks and gauges, \eqref{eq:original-15} for the spatial
readings still uses just the appropriate field/tree data in the resulting coordinate, not any
additional information at the normalization positions.

\begin{lemma}[Field reconstruction]\label{conf:field-reconstruction}
In the comparison classes above, the limiting critical area measure and signed colored \(O\) data
determine \(h\) and \(J\), respectively.
\end{lemma}
\begin{proof}

For \(J\) this is the whole-plane reconstruction of Section~\ref{sec:walls}, including the parity
convention. For the real field we use the log-mass reconstruction strategy
of Berestycki--Sheffield--Sun \citep[Section~2]{BSSReconstruction}; critical
reconstruction under Gaussian covariance regularity hypotheses was established
by Vihko \citep[Theorem~A]{VihkoReconstruction}.  We give the local argument
needed for our critical-area normalization and comparison classes. In a patch of an ordinary pinned log GFF
with covariance \(\log(1/|v-w|)+\mathrm{const}\), at an inset \(x\) subtract its circle average
\(h_\eta(x)\) and zoom by \(\eta\). On the unit open disk the law of the result \(Y_{\eta,x}\)
(scalar pullback here) is fixed with pure log covariance. For two distinct inset points these
distributions jointly converge to independent copies by Gaussian covariances. With \(\mu=\mu_h^{\rm
crit}\), field and scale covariance give
\[
 \tfrac12\log\mu(B(x,\eta/4))+2\log(1/\eta)
                =h_\eta(x)+\tfrac12\log\mu_{Y_{\eta,x}}^{\rm crit}(B(0,1/4)).
\]
Write \(M\) for the unit-disk mass in the last term. Its logarithm
has finite moments of all orders. For the upper tail, local chaos
convergence and Fatou's lemma transfer the positive fractional moment
of Lemma~\ref{cont:fractional-moment} to the critical Dirichlet mass.
A Markov decomposition on a larger disk adds a harmonic field whose
supremum on \(B(0,1/4)\) has Gaussian tails. H\"older's inequality,
with a smaller positive exponent if necessary, therefore gives
\(\E M^s<\infty\) for some \(s>0\), and hence
\[
                         \Prob(\log M>m)\le C e^{-sm}.
\]

For the lower tail, pack \(N_m\ge c_0e^{2cm}\) disjoint balls of
radius comparable to \(e^{-cm}\) inside \(B(0,1/4)\), where
\(c>0\) will be small. Decompose into independent zero-Dirichlet
fields inside them and harmonic parts. Gaussian harmonic estimates and
the grid union bound give an event of probability at least
\(1-Ce^{-cm}\) on which all the harmonic parts on the half-balls
are at least \(-C_0cm\). Critical scaling then gives, on that event,
\[
 M\ge e^{-(4+2C_0)cm}\max_{1\le j\le N_m} Z_j
\]
after adjusting fixed constants, where the \(Z_j\) are independent
copies of a positive unit Dirichlet inner-ball mass. Choose
\(z_0>0\) with \(p_0=\Prob(Z_j\ge z_0)>0\), and then choose
\(c\) so that \((4+2C_0)c<1\). For all sufficiently large \(m\),
\[
 \Prob(\log M<-m)
 \le Ce^{-cm}+(1-p_0)^{N_m}.
\]
This proves the asserted log-moment bounds.

By fixed marginal laws and measurability, the two log masses at distinct
points converge jointly to independent copies, by Lusin approximation.
The log-moment bounds make their covariance tend to zero. Put
\[
 m_0=\tfrac12\E\log M,\qquad
 R_\eta(x)=\tfrac12\log\mu(B(x,\eta/4))
                         +2\log(1/\eta)-m_0.
\]
The log-mass identity shows that \(R_\eta-h_\eta\) has mean zero.
Integrating its covariance against any two smooth test functions on the
inset patch and using the uniform second-moment bound gives convergence
to zero in \(L^2\). Since \(h_\eta\to h\) in tested \(L^2\), the
measure-based functions \(R_\eta\) reconstruct \(h\). A subsequence
and a countable determining family of tests give almost-sure measurable
reconstruction.

Local actual-height absolute continuity off ends transfers this in ordinary spherical charts of the
unpinned experiments (express \(h\) in planar area coordinates). One need not know end locations for
determination: two conditional copies given the mass agree throughout the complement of both finite
end sets by such patches, hence agree everywhere by \(H^{-1}\). Indeed a distribution supported on finitely many
planar points is a finite linear combination of delta masses and their derivatives, and none of these
nonzero distributions belongs to \(H^{-1}\) in dimension two.

\end{proof}

\begin{proposition}[Removal of the auxiliary sampling choices]\label{conf:source-elimination}
Given \(h,J\), the conditional laws of the regularized true profile arrays are unchanged by
revealing the central loop, associated end choices and phase, and the entire canonical source
family. Conditional profile kernels for disjoint strictly buffered observations multiply given
\(h,J\).
\end{proposition}
\begin{proof}
We use a conditional-sampling principle for the true profile array
\(Z\). If an auxiliary choice \(C\), sampled after the finite map,
has in the limit a conditional law given \((h,J,Z)\) that depends
only on \((h,J)\), then \(C\) and \(Z\) are conditionally
independent given \((h,J)\). We verify this first for the central
choices and then, round by round, for the source family. The spatial
product kernels are used with the sampled field and embedding retained
in the conditioning; they are not obtained by averaging a shared
embedding out of the conditioning variables.

\smallskip\noindent\emph{The central loop, ends, and phase.}\par

Use \(b_n\)-type unconditioned bases, with \(B_n\) nearest the desired units, localized by
first-perimeter windows and further cuts with both central side fractions positive and total area in
compact positive ranges. These exhaust the arbitrary-area Palm measure class and agree on overlaps
up to input weights. Conditional on the finite map and its orientations beforehand, the joint mass
is exactly loop counting with the indicated restrictions, up to factors depending only on the full
map: this is the same \(4/(k+l)\) ring calculation in \eqref{eq:original-3} with the total not fixed
(on maps of \(m\) edges the change from the root contributes a factor common at that \(m\)).

We can take two additional end samples by area, requiring separation, and restore any equivalent
positive counting weights by limiting changes of density. These are not required to be the three
conformal normalization samples. Interior continuity-window cutoffs exhaust.

The entire counting and end-sampling conditional rule now passes to the limit conditional on map
outputs (including true profile arrays): diffuse weak area convergence excludes small loops under
the side cutoff and samples their positions, and all relevant loops, their perimeter labels and side
measures converge, with null boundary area. Equivalently, under central counting within such
restrictions one averages uniformly over all qualifying loops given the map, then end samples on
opposite sides by their area laws (or biases by the separation probability). This conditional
sampler converges and depends only on \(h,J\) in the limit. The common typical central phase has the
same assertion by length-label convergence.

All available loop choices and opposite end samplings are covered by increasing these restrictions.
Conditional profile versions on different windows come from the same extractions/formulas, so this
proves independence from these choices given \(h,J\), in the full supported class (each two choices
occur eventually in overlapping experiments). Independent normalization marks before the gauge
change used map area; no loop selection or source operation is used to generate that initial true
profile array.

\smallskip\noindent\emph{A finite-round survey of caller births.}\par

We detail similarly the source assertion, since disk identities might be reached adaptively. Run a
countable survey in rounds starting at the two central sides. At a complete caller birth being
entered, draw its source freshly before testing there. In previously entered callers test paths to a
dense sequence of ordinary length targets, on compact surviving prefixes, using preassigned loop
orientations. Other original-target branches use the same caller start and the side-dependent pinch
switch rule, so testing whole individual paths here requires no new independent source along an
inherited-target portion. List macroscopic full births discovered, including a negative at which the
inherited interval is empty, and append new callers to the entered list, drawing sources there for
the further tests.

Repeatedly increase target lists, time cutoffs and resolution of jump sizes and of minimum positive
surviving side length, also in old callers. Target seeds and continuity cutoffs can be diffuse fresh
choices; every round uses finitely many tests/births almost surely. For example use increasing clock
bounds and shrinking lower bounds before the first small length, taking only jumps with strict
compactly surviving pre/post comparisons; order listed births by the path test and chronology,
skipping duplicates.

At each fixed round the decisions/callers and slot measures at their births are matched stably along
good subextractions by the \(J_1\), fork/empty-interval, shared-prefix and length-label spatial
comparisons. Bounds can avoid discontinuities by the fresh cutoffs. Shared births met in finitely
many paths are along their common prefixes with strict original target eligibility.

The countable survey enters every complete macroscopic caller: every next birth is on an
ordinary-target open interval in its parent with a compact surviving jump observation, and depth can
be taken successively. Use a fair schedule on all previously entered callers.

All these finite tests can be applied to the same primary source assignments on the lattice. Only
actually needed independent slots have to be drawn as of a query (extra not-yet-examined calls
deferred); one can skip a proposed fresh draw if it is not a valid new birth, checking this by the
prior histories. No primary absolute priority among targets is imposed.

\smallskip\noindent\emph{Passage of the conditional samplers to the limit.}\par

Thus conditionally on the map, marks and the survey past, every listed fresh slot has its uniform
distribution at that birth. Passing this equality against bounded continuous tests to the limit uses
the matched uniform boundary-slot measure convergence in the true conformal coordinates for that
birth (negative as well as positive). Failures/cutoff tails in any fixed round can first be
truncated.

Hence given the limiting map outputs, central choices and past, the new source law is the
corresponding normalized disk length law. Those outputs include enough to give \(h,J\) and the true
regularized profiles by the preceding determination; survey seeds may likewise be drawn
progressively independently of them.

Given fields, earlier choices and tested seeds, the limiting history deciding the next draw uses
exactly the canonical level paths and their physical lengths/clocks determined above, without a
further profile variable. Iteration and then round exhaustion proves the assertion for the whole
source family, including with its intrinsic caller identities in that exploration.

All these independence assertions concern profile readings which could themselves have been
subextracted from countable true map tests before choosing loops/sources; no independence of a
microscopic exploration array from the map is being claimed.

Consequently conditional profile kernels on disjoint strictly buffered observations multiply given
just \(h,J\). Indeed they did so given the whole input by \eqref{eq:original-15}, the component
kernels are now unchanged after dropping the source/loop/phase/end information, and the embedded
primary path input itself is given by the fields and these choices. All-candidate bounds over field
extensions persist.

\end{proof}

\begin{lemma}[Locality and constant-on-buffer comparisons]\label{conf:field-kernel-locality}
A protected profile test contained in the bounded-side Jordan disk of a whole \(O\) loop has its
field kernel determined by any field patch containing that loop and disk. More generally, if two
typical supported inputs have the same \(J\) and their real fields differ by a deterministic bounded
smooth function which vanishes near the test, then the two kernels agree. If that function equals a
constant \(b\) near the test, its kernel agrees there with that of \(h+b\). These assertions concern
common deterministic cost units and do not assert Weyl scaling of graph cost.
\end{lemma}
\begin{proof}
\smallskip\noindent\emph{Localization inside a whole loop.}\par

In a planar patch containing a whole \(O\) loop whose bounded-side Jordan disk \(D\) strictly
contains the protected test, the field kernel for that test is local up to this field patch. One can
use the isolated disk calculation in the germ proof, rerooting if necessary so that a central choice
exterior to \(D\) and ends there are used and \(D\) is a positive birth. Such choices have positive
conditional support by nested CLE and positive area. Bounded planar tests have this determination
with probability arbitrarily near one by increasing a finite surrounding field buffer, since
bounded-side \(O\) disks enclose each given compact (use nesting toward infinity).

\smallskip\noindent\emph{Smooth field changes and simultaneous observations.}\par

For \(h'\) related to \(h\) by a deterministic bounded smooth Weyl change in typical supported
inputs with the same \(J\), suppose the two real fields agree in a strict neighborhood of one
observation. Their kernels there are identical. This also compares \(h+f\) to \(h+b\) when \(f\) is
constant \(b\) there. Use the exact same-chart comparison in Lemma~\ref{kern:field-comparison}.

The common central loop/order and the needed caller sources can be shared when making it: globally
large callers (and required ancestors) are finite, with mutually absolutely continuous
length-sampling laws; smaller ones required throughout whole matching neighborhoods have the
identical internal laws, and unrelated subtrees entering no retained observation need not be shared.
Alternatively couple all callers above a fixed small diameter bound and all smaller ones meeting a
slightly inset comparison region. Boundary measures change by a continuous positive weight. By the
source/loop independence just proved, this establishes equality of the field kernels themselves
almost surely without a likelihood price per crossing caller.

When several disjoint patches are considered, use conditional product also to sample auxiliary
kernels there independently (including at specified constant shifts); one can thus keep designated
inner constant-comparison and exterior unchanged readings in the modified law together. All such
statements use joint extractions with fixed deterministic units; they are not Weyl scaling for graph
costs.

\end{proof}

\subsection{Full-field Cameron--Martin comparisons}

\begin{proposition}\label{conf:cm-comparison}
Compactly supported deterministic Cameron--Martin changes in an ordinary observation patch preserve
the full-field arbitrary-area measure class. Quantitative comparisons for any finite list of smooth
recipes in similar shells away from the ends can be made with bounded likelihood on events of
arbitrarily high marginal probability, using ordinary local Gaussian pairings and energies. The
exterior can be retained in these comparisons.
\end{proposition}
\begin{proof}
\smallskip\noindent\emph{A bridge-and-tail description of the measure class.}\par

Randomize the cylinder maximum position and height with positive densities, with two ends and
full-support spherical re-embedding; positive global Palm and mark weights do not change this
arbitrary-area class.

For cutoffs \(m,K>0\) use comparison laws with full-density endpoint heights at \(\pm m\) below
\(K\), a Brownian bridge of circle means inside \([-m,m]\), and independent tails of \(K-{\rm
BES}_3\) from those endpoints; use the same independent lateral field. Each such law is dominated in
measure class by the arbitrary-height/position maximum law, which conversely is dominated by
countable mixtures of them.

Indeed in a longer interval \([-T,T]\) localize the maximum construction to maximum attained there
with height in \((-B,B)\). Its interior path then uses the Brownian bridge class with those endpoint
heights of positive densities and the given maximum/time (both have positive joint conditional
density on the admissible range); the two independent Bessel pieces are precisely the Brownian
maximum decomposition bridges (Denisov decomposition, or the killed Brownian bridge density
differentiated in the maximum).

The tails are killed Brownian height paths conditioned by the linear
avoid-the-maximum transform. To compare barriers, start at \(x<M_-<M_+\).
For the height process \(M_+-{\rm BES}_3\), the probability of staying
below \(M_-\) forever is
\[
                         \frac{M_--x}{M_+-x}>0.
\]
Conditioning on this event gives exactly the \(M_--{\rm BES}_3\) tail
law started at \(x\). Indeed the two Doob transforms have harmonic
functions \(M_+-x\) and \(M_--x\), whose ratio tends to one at minus
infinity. Thus a lower-barrier tail is an absolutely continuous
restriction of a higher-barrier tail, with the displayed normalizing
factor. Apply this with the actual maximum as the lower barrier and,
for example, \(B+1\) as the upper barrier. Thus on these restrictions the maximum law has the same
class as Brownian interior/full-density endpoints and higher-barrier Bessel tails, restricted to
inside maximum dominating both tails and in the stated height window. This proves one domination.

Conversely any of the comparison laws clipped to these maximum/time restrictions at larger \(T,B\)
(\(B>K\)) is dominated by the same class: finite segments of its stopped-height tails are ordinary
killed Brownian paths with endpoint tilts and the residual tails use a lower barrier. Exhaust using
convergence of tails to minus infinity. Apply these comparisons also with full positive re-embedding
densities, conditioning on such end-coordinate parameters when convenient.

\smallskip\noindent\emph{Quantitative changes within the slab.}\par

On an observation patch inside such a Brownian slab, compactly supported deterministic
Cameron–Martin changes therefore preserve the full-field measure class.

For quantitative changes one can actually use the bridge comparison laws with a large enough slab,
clipping to inputs of interest beforehand. The likelihood for a smooth compact shift there only uses
the ordinary Dirichlet Gaussian pairing and energy of the shift (conditional Brownian linear mean
gives no term), retaining the exterior, and the independent lateral part with the same Dirichlet
normalization. Conformal chart expression just includes the known background transform.

Thus for finite recipe lists in small similar shells off ends one can impose bounded likelihood
tests of arbitrarily high marginal probability by Gaussian local cutoffs in their aligned shell
coordinates. The comparison is made within the chosen Brownian-slab law,
with its exterior retained. It does not require a uniform bound on the
original sphere's global area-tilt likelihood at each shrinking shell.

\end{proof}
